% =====================================================================================
%  Chapter 7 — Multi-Robot and Scaling Results
%  Thesis: Comparison of Control Strategies for Interaction-Force Aware Multirotor Teams
%
%  DRAFT SKELETON, 27 September 2026.
%  Sources: docs/30_Results_Chapters_Skeleton.md (Chapter 7), docs/36_Four_Five_Robot_Formation_Sketch.md,
%           Chapter~\ref{ch:setup} (three-robot scenarios B1--B3).
%
%  HONESTY NOTE. Three-robot C.4 data may be partial or absent; the fallback paragraph in
%  Section~\ref{sec:res3-limitation} is written now as the contingency narrative. n$\geq$4 content
%  is explicitly future work, not a claim of results.
%
%  Citation baseline: 2026-08-27 (see docs/thesis/CITATION_BASELINE.md).
% =====================================================================================

\chapter{Multi-Robot and Scaling Results}
\label{ch:results3}

RQ3 (Section~\ref{sec:intro-rq}) asks for the transfer penalty when a summed per-neighbour
residual model trained on two-vehicle data is applied to three-vehicle formations. This chapter
reports that penalty on the B-series scenarios if three-robot comparison flights are flown, and
states how an extension to four or five vehicles would be structured without treating that extension
as part of the primary thesis gate.

% -------------------------------------------------------------------------------------
\section{Three-robot comparison campaign}
\label{sec:res3-campaign}

Three-robot flights require the third brushless vehicle in the roster and scenarios B1, B2 and/or
B3 from the frozen library (I-stack, V-stack, and merge geometries). The metric pipeline is
intentionally identical to Chapter~\ref{ch:results2}: phase metrics, separation error, effort
proxies, and residual logging under Strategy~0 versus compensating strategies where implemented
for $N=3$.

\todo[color=red!25]{PLACEHOLDER: Table --- B1/B2/B3 (subset as flown) $\times$ Strategy~0--3;
same column headers as Table~\ref{tab:res2-tracking}; $n$ repeats per cell.}

% -------------------------------------------------------------------------------------
\section{Transfer penalty for the summed predictor}
\label{sec:res3-transfer}

RQ3 is not a test of whether aerodynamic interactions superpose; the literature already reports that
they do not. The engineering question is how large the penalty is when the same onboard predictor
trained on two-robot C.1 data is evaluated on three-robot relative configurations, and whether that
penalty is small enough to keep the summed architecture.

\todo[color=red!25]{PLACEHOLDER: Figure/table --- open-loop prediction error on three-robot logs
vs two-robot training distribution; per-neighbour vs summed prediction; scenario-resolved breakdown.
Source: C.1/C.2 artifacts and merged three-robot uSD.}

\todo[color=red!25]{PLACEHOLDER: Closed-loop block --- if Strategy~2/3 flown on B-series, compare
tracking/residual metrics to Strategy~0 on the same scenarios.}

% -------------------------------------------------------------------------------------
\section{Extension to four and five vehicles}
\label{sec:res3-n4sketch}

A sketch of scenarios with four or more vehicles describes how such formations would stress
superposition and asymmetric wash without replacing the frozen scenario library. The sketch is a
template only: its geometries must pass the pre-flight check and the supervisor's approval before any
simulation or flight. No results from it are claimed here.

The primary thesis comparison remains two- and three-robot; n$\geq$4 is documented so a future
campaign can reuse the same logging, merge and analysis tools.

% -------------------------------------------------------------------------------------
\section{Limitation if three-robot data are incomplete}
\label{sec:res3-limitation}

\paragraph{If the third vehicle arrives too late.}
RQ3 needs flights with three vehicles, because the effect it asks about, the overlap of two wakes,
does not exist in two-vehicle data. If the third vehicle is not available in time, RQ3 is not
answered on hardware, and the thesis says so. Two-vehicle logs can only check how the network behaves
on one neighbour at a time. They say nothing about the summed prediction for two neighbours.

The limitation is stated plainly: without repeated three-robot compensated flights, the thesis
cannot rank Strategies~1--3 at $N=3$ with the same confidence as at $N=2$; the summed-architecture
decision rests on the magnitude of open-loop penalty and simulation sanity checks, and generalisation
to larger teams remains future work (Section~\ref{sec:conclusion-future}).

% -------------------------------------------------------------------------------------
\section{Chapter summary}
\label{sec:res3-summary}

\todo[color=red!25]{PLACEHOLDER: One paragraph summarising RQ3 answer once metrics exist --- penalty
size, whether summed predictor remains defensible, link back to Section~\ref{sec:res2-summary} for
two-robot vs three-robot consistency.}
